\begin{textsample}{NEG Dim 4 – 2021\_c4 – Score 1.00 – 2021\_c4/enade\_2021\_filo\_baa597e9cf.txt}  \label{ex:f4_neg_224}
CAPÍTULO 4 PERCEPÇÃO DA PROVA Nas análises feitas neste capítulo, constam as percepções dos concluintes da Área de Filosofia ( Licenciatura e Bacharelado ) sobre a prova aplicada no Enade/2021.
Essas percepções foram mensuradas por meio de nove questões que avaliaram desde o grau de dificuldade da prova até o tempo gasto para concluí -la.
As percepções sobre a prova foram relacionadas ao desempenho dos estudantes, à Grande Região de funcionamento do curso, à Categoria Administrativa e ao tipo de Organização Acadêmica da IES.
O questionário de percepção da prova encontra -se ao final dos Anexos VIIIA ( Licenciatura ) e VIIIB ( Bacharelado ), onde está a reprodução do exame.
O desempenho dos estudantes foi classificado em quatro quartos.
Para tanto, esse desempenho foi ordenado de forma ascendente.
O percentil 25, P25, também conhecido como primeiro quartil, é a nota de desempenho que deixa um quarto ( 25 \% ) dos valores observados abaixo e três quartos acima.
Na Figura 1, apresenta -se uma ilustração desse conceito.
O quarto inferior de desempenho é composto pelas notas abaixo do primeiro quartil.
Já o percentil 75, P75, também conhecido como terceiro quartil, é o valor para o qual há três quartos ( 75 \% ) dos dados abaixo e um quarto acima dele.
O quarto superior de desempenho é composto pelas notas iguais ou acima do terceiro quartil.
O percentil 50, P50, também conhecido como mediana, é o valor que divide as notas em dois conjuntos de igual tamanho.
O segundo quarto inclui os valores entre o primeiro quartil ( P25 ) e a mediana.
O terceiro quarto contém os valores entre a mediana ( P50 ) e o terceiro quartil ( P75 ).
Vale ressaltar que percentis, quartis e medianas são pontos que, não obrigatoriamente, pertencem ao conjunto original de dados, ao passo que os quartos são subconjuntos dos dados originais.
A seguir, serão apresentados gráficos com resultados selecionados, relativos às nove questões avaliadas por grupos de estudantes.
Nas barras dos gráficos, apresenta -se o percentual de alunos que assinalaram uma das opções ou a soma das porcentagens daqueles que assinalaram duas ( ou três ) delas.
Por exemplo, para as questões 1 e 2, nos gráficos, é apresentada a porcentagem total de participantes que assinalaram as opções ( D ) Difícil e ( E ) Muito difícil.
Em cada barra, foram assinalados, também, os extremos do intervalo de confiança de 95 \% como linhas verticais unidas por uma linha horizontal na forma da letra H maiúscula.
O estimador de um parâmetro com um certo nível de confiança ( por exemplo, 95,0 \% ) deve conter o parâmetro no intervalo de confiança em 95 \% das vezes.
Na comparação entre os estimadores dos parâmetros de duas classes de uma dada categoria ( por exemplo : Norte e Nordeste nas Grandes Regiões, de primeiro e último quarto dentro de desempenho, de IES Públicas e Privada, ou de Universidades e Faculdades ), associados aos seus respectivos intervalos de confiança, diz -se que não há diferença estatisticamente significativa entre os parâmetros das duas categorias quando há interseção entre os intervalos de confiança, e que há diferença, se os intervalos de confiança são disjuntos ( para maiores informações vide Glossário ).
Nas Tabelas do Anexo IIA ( Licenciatura ) e do Anexo IIB ( Bacharelado ), são apresentados os valores absolutos e a distribuição percentual das alternativas válidas das nove questões, segundo o mesmo recorte de desempenho dos alunos e da Grande Região de funcionamento do curso, de Categoria Administrativa e do tipo de Organização Acadêmica da IES. 4.1 GRAU DE DIFICULDADE DA PROVA 4.1.1 COMPONENTE DE FORMAÇÃO GERAL Ao avaliarem “ Qual o grau de dificuldade desta prova na parte de Formação Geral? ” ( Questão 1 ), 29,0 \% do grupo de inscritos e presentes de Licenciatura optaram pelas alternativas Difícil ou Muito difícil.
Entretanto, para mais da metade dos estudantes ( 57,9 \% ), o Componente de Formação Geral da prova foi avaliado como de grau de dificuldade Médio ( Gráfico 4.1a, Gráfico 4.2a e, no Anexo IIA, a Tabela II.1a ).
O percentual de estudantes de Licenciatura que consideraram a prova Difícil ou Muito difícil foi maior na região Norte, onde a proporção foi de 36,4 \%, enquanto a de menor proporção foi na Nordeste, com 27,5 \%.
No Gráfico 4.1a, é possível observar que as diferenças entre as regiões não são estatisticamente significativas.
Nas Grandes Regiões, a proporção de presentes à prova que avaliaram o Componente de Formação Geral como de grau de dificuldade Médio esteve entre 51,7 \%, na região Norte, e 61,4 \%, na região Nordeste. entre as regiões não são estatisticamente significativas.
Nas Grandes Regiões, a proporção de presentes à prova que avaliaram o Componente de Formação Geral como de grau de dificuldade Médio esteve entre 51,7 \%, na região Norte, e 61,4 \%, na região Nordeste.
Ao avaliarem “ Qual o grau de dificuldade desta prova na parte de Formação Geral? ” ( Questão 1 ), 27,3 \% do grupo de inscritos e presentes de Bacharelado optaram pelas alternativas Difícil ou Muito difícil.
Entretanto, para mais da metade dos estudantes de Bacharelado ( 56,6 \% ), o Componente de Formação Geral da prova foi considerado como de grau de dificuldade Médio ( Gráfico 4.1b, Gráfico 4.2b e, no Anexo IIB, a Tabela II.1b ).
O percentual de estudantes de Bacharelado que consideraram a prova como Difícil ou Muito difícil foi maior na região Norte, onde a proporção foi de 43,2 \%, enquanto a de menor incidência foi na Sudeste, com 24,1 \%.
No Gráfico 4.1b, é possível observar que as diferenças entre as regiões não são estatisticamente significativas.
Nas Grandes Regiões, a proporção de presentes à prova que consideraram o Componente de Formação Geral como sendo de grau de dificuldade Médio esteve entre 35,1 \%, na região Norte, e 60,8 \%, na região Nordeste.
O percentual de alunos de Licenciatura que consideraram o componente de Formação Geral da prova Difícil ou Muito difícil foi decrescente em função dos quartos de desempenho : 40,3 \%, no primeiro quarto, e 18,6 \%, no quarto, o grupo de melhor desempenho na prova.
Nos quartos de desempenho intermediários, a proporção de alunos que consideraram a prova Difícil ou Muito difícil foi de 32,0 \%, no segundo quarto, e de 25,6 \%, no terceiro quarto.
As diferenças entre os quartos de desempenho são estatisticamente significativas.
Para todos os quartos de desempenho, a alternativa modal a essa pergunta foi a Médio, com 50,8 \% e 62,3 \% dos respondentes no primeiro e terceiro quartos, respectivamente.
O percentual de alunos de Bacharelado que consideraram o componente de Formação Geral da prova Difícil ou Muito difícil foi decrescente em função dos quartos de desempenho : 37,3 \%, no primeiro quarto, e 18,7 \%, no quarto, grupo de melhor desempenho na prova.
Nos quartos de desempenho intermediários, a proporção de alunos que consideraram a prova Difícil ou Muito difícil foi de 30,1 \%, no segundo quarto, e de 23,5 \%, no terceiro quarto.
As diferenças entre o quarto inferior e os dois quartos superiores de desempenho são estatisticamente significativas, assim como entre o segundo quarto e o quarto superior.
Para todos os quartos de desempenho, a alternativa modal a essa pergunta foi a Médio, com 51,9 \% e 59,3 \% dos respondentes no primeiro e no segundo quartos, respectivamente.
O percentual de estudantes de Licenciatura que consideraram o componente de Formação Geral da prova Difícil ou Muito difícil foi maior entre os estudantes de Centros Universitários, alcançando 31,3 \%, enquanto entre os estudantes das Faculdades, a Organização Acadêmica com a menor incidência, alcançou 25,4 \%.
No Gráfico 4.3a, é possível observar que as diferenças entre as Organizações Acadêmicas não são estatisticamente significativas.
Nas Organizações Acadêmicas, a proporção de presentes de Licenciatura à prova que avaliaram o Componente de Formação Geral como de grau de dificuldade Médio esteve entre 57,7 \%, nas Faculdades, e 58,0 \% nos Centros Universitários ( ver também a Tabela II.2a, no Anexo IIA ).
O percentual de estudantes de Bacharelado que consideraram o componente de Formação Geral da prova como Difícil ou Muito difícil foi maior entre os estudantes de Centros Universitários, uma proporção de 38,1 \%, enquanto nas Universidades, a Organização Acadêmica com a menor incidência, a proporção atingiu 22,8 \%.
No Gráfico 4.3b, é possível observar que apenas a diferença entre Faculdades ( 33,6 \% ) e Centros Universitários não é estatisticamente significativa.
Nas Organizações Acadêmicas, a proporção de presentes de Bacharelado à prova que consideraram o Componente de Formação Geral como sendo de grau de dificuldade Médio esteve entre 50,7 \%, nas Faculdades, e 59,2 \%, nas Universidades ( ver também a Tabela II.2b, no Anexo IIB ).
O percentual de estudantes de Licenciatura que consideraram o componente de Formação Geral da prova Difícil ou Muito difícil foi maior entre os estudantes de cursos de IES Privadas ( 31,3 \% ), sem diferença estatisticamente significativa da proporção registrada em IES Públicas ( 27,1 \% ).
Já a proporção de presentes de Licenciatura à prova que avaliaram o Componente de Formação Geral como de grau de dificuldade Médio foi de 59,2 \%, entre os estudantes de IES Públicas, e de 56,4 \%, entre os de IES Privadas ( ver Gráfico 4.4a ).
O percentual de estudantes de Bacharelado que consideraram o componente de Formação Geral da prova como Difícil ou Muito difícil foi maior entre os estudantes de cursos de IES Privadas ( 31,2 \% ), uma diferença estatisticamente significativa da proporção registrada em IES Públicas ( 21,6 \% ).
Já a proporção de presentes de Bacharelado à prova que consideraram o Componente de Formação Geral como sendo de grau de dificuldade Médio foi de 58,2 \%, entre os estudantes de IES Públicas, e de 55,5 \%, entre os de IES Privadas ( ver Gráfico 4.4b ). 4.1.2 COMPONENTE DE CONHECIMENTO ESPECÍFICO Ao responderem à Questão 2 – “ Qual o grau de dificuldade desta prova na parte de Componente Específico? ” – 39,6 \% do grupo de estudantes de Licenciatura classificaram -na como Difícil ou Muito difícil.
Além disso, o Componente modal de Conhecimento Específico da prova foi considerado de grau de dificuldade Médio por 52,6 \% dos alunos ( Gráfico 4.5a, Gráfico 4.6a, e, no Anexo IIA, a Tabela II.3a ).
A análise das respostas dos estudantes de Licenciatura quanto ao grau de dificuldade do Componente de Conhecimento Específico da prova, agregado por Grande Região, mostra diferenças estatisticamente significativas entre proporções de alunos que a avaliaram como Difícil ou Muito difícil.
Tais diferenças são observadas entre a região Sul ( 45,1 \% ) e as regiões Nordeste ( 35,6 \% ) e Sudeste ( 37,2 \% ).
O percentual de alunos que classificaram o grau de dificuldade como Médio, no Componente de Conhecimento Específico, variou de 46,3 \% a 56,6 \%, nas regiões Sul e Nordeste, respectivamente.
Ao responderem à Questão 2 – “ Qual o grau de dificuldade desta prova na parte de Componente Específico? ” – 40,7 \% do grupo de estudantes de Bacharelado classificaram -na como Difícil ou Muito difícil.
Além disso, o Componente modal de Conhecimento Específico da prova foi considerado como de grau de dificuldade Médio por 52,0 \% dos alunos ( Gráfico 4.5b, Gráfico 4.6b, e, no Anexo IIB, a Tabela II.3b ).
A análise das respostas dos estudantes de Bacharelado quanto ao grau de dificuldade do Componente de Conhecimento Específico da prova, agregado por Grande Região, mostra uma diferença estatisticamente significativa entre proporções de alunos que a avaliaram como Difícil ou Muito difícil.
Tal diferença é observada apenas entre a região Sudeste ( 36,8 \% ) e a região Sul ( 47,1 \% ).
O percentual de alunos que classificaram o grau de dificuldade como Médio, no Componente de Conhecimento Específico, variou de 29,7 \% a 55,5 \%, nas regiões Norte e Sudeste, respectivamente.
Considerando -se a avaliação da dificuldade das questões do Componente de Conhecimento Específico da prova de acordo com o desempenho dos estudantes de Licenciatura, observa -se uma tendência decrescente dos que consideraram a parte específica Difícil ou Muito difícil, variando de 29,1 \% ( último quarto ) a 46,6 \% ( primeiro quarto ).
As demais proporções foram de 42,1 \% e 40,7 \%, respectivamente, no segundo e terceiro quartos.
Observa -se que há diferença estatisticamente significativa dos resultados entre o quarto superior de desempenho e os demais.
Já os que responderam que o grau de dificuldade das questões do Componente de Conhecimento Específico da prova foi Médio atingiram percentagens que variaram de 47,1 \%, no primeiro quarto, a 59,4 \%, no último.
Considerando -se a avaliação da dificuldade das questões do Componente de Conhecimento Específico da prova, de acordo com o desempenho dos estudantes de Bacharelado, observa -se que não há diferença estatisticamente significativa dos resultados entre os quartos de desempenho.
A proporção dos que classificaram a parte específica como Difícil ou Muito difícil variou de 34,3 \% ( último quarto ) a 43,4 \% ( terceiro quarto ).
As demais proporções dessa resposta foram de 42,0 \% e de 43,1 \%, nos quartos inferior e segundo, respectivamente.
Já a proporção dos que responderam que o grau de dificuldade das questões do Componente de Conhecimento Específico da prova foi Médio variou de 48,4 \%, no terceiro quarto, a 57,8 \%, no último.
Considerando -se a avaliação da dificuldade das questões do Componente de Conhecimento Específico da prova, o percentual de estudantes de Licenciatura que as consideraram Difícil ou Muito difícil foi maior entre os estudantes de Centros Universitários, atingindo 45,9 \% ; já entre os de Universidades, a Organização Acadêmica com a menor incidência, essa proporção foi de 37,3 \%.
No Gráfico 4.7a, é possível observar que apenas a diferença entre Universidades e Centros Universitários é estatisticamente significativa.
Nas Organizações Acadêmicas, a proporção de presentes de Licenciatura à prova que consideraram o Componente de Conhecimento Específico de grau de dificuldade Médio esteve entre 47,1 \%, nos Centros Universitários, e 54,7 \%, nas Universidades ( ver também Tabela II.4a, no Anexo IIA ).
Considerando -se a avaliação da dificuldade das questões do Componente de Conhecimento Específico da prova, o percentual de estudantes de Bacharelado que as consideraram como Difícil ou Muito difícil foi maior entre os estudantes de Centros Universitários, uma proporção de 51,3 \%, enquanto entre os de Universidades, a Organização Acadêmica com a menor incidência, o índice atingiu 37,1 \%.
No Gráfico 4.7b, é possível observar que apenas a diferenças entre Universidades e Centros Universitários é estatisticamente significativa.
Nas Organizações Acadêmicas, a proporção de presentes de Bacharelado à prova que consideraram o Componente de Formação Geral como sendo de grau de dificuldade Médio esteve entre 45,4 \%, nos Centros Universitários, e 54,6 \%, nas Universidades ( ver também Tabela II.4b, no Anexo IIB ).
O percentual de estudantes de Licenciatura que consideraram o Componente Específico da prova Difícil ou Muito difícil foi maior entre os estudantes de cursos de IES Privadas ( 43,7 \% ), com diferença estatisticamente significativa da proporção registrada em IES Públicas ( 35,9 \% ).
Já a proporção de presentes de Licenciatura à prova que avaliaram esse componente como de grau de dificuldade Médio foi de 55,8 \%, entre os estudantes de IES Públicas, e de 48,9 \%, entre os de IES Privadas ( ver Gráfico 4.8a e Tabela II.4a, no Anexo IIA ).
O percentual de estudantes de Bacharelado que consideraram o Componente Específico da prova como Difícil ou Muito difícil foi maior entre os estudantes de cursos de IES Privadas ( 43,4 \% ), sem uma diferença estatisticamente significativa da proporção registrada em IES Públicas ( 36,5 \% ).
Já a proporção de presentes de Bacharelado à prova que consideraram esse componente como sendo de grau de dificuldade Médio foi de 53,9 \% entre os de IES Públicas, e de 50,7 \%, entre os de IES Privadas ( ver Gráfico 4.8b e Tabela II.4b, no Anexo IIB ). 4.2 EXTENSÃO DA PROVA EM RELAÇÃO AO TEMPO TOTAL Indagados quanto à extensão da prova, em relação ao tempo total oferecido para a sua resolução ( Questão 3 ), os estudantes de Licenciatura apontaram, com maior incidência, a alternativa que considerava a extensão Adequadas, em todas as agregações consideradas ( Gráfico 4.9a, Gráfico 4.10a, e, no Anexo IIA, a Tabela II.5a ).
O percentual de alunos de Licenciatura que responderam ser a extensão da prova Adequada foi de 46,7 \%.
Já 43,2 \% dos inscritos presentes consideraram que a prova foi Longa ou Muito longa, e 10,1 \% a avaliaram como Curta ou Muito curta.
Entre as Grandes Regiões, os que avaliaram a prova como Longa ou Muito longa, em relação ao tempo total destinado à sua resolução, atingiram uma porcentagem que variou de 33,9 \%, na região Norte, até 43,9 \%, na região Sudeste.
Não há diferença estatisticamente significativa entre as regiões.
Indagados quanto à extensão da prova, em relação ao tempo total oferecido para a sua resolução ( Questão 3 ), os estudantes de Bacharelado apontaram, com maior incidência, a alternativa que considerava a extensão Adequada, em todas as agregações consideradas ( Gráfico 4.9b, Gráfico 4.10b, e, no Anexo IIB, a Tabela II.5b ).
O percentual de alunos de Bacharelado que responderam ser a extensão da prova Adequada foi de 44,2 \%.
Já 46,1 \% dos inscritos presentes consideraram que a prova foi Longa ou Muito longa, e 9,8 \% a avaliaram como Curta ou Muito curta.
Entre as Grandes Regiões, a proporção daqueles que avaliaram a prova como Longa ou Muito longa, em relação ao tempo total destinado à sua resolução, variou de 41,0 \%, na região Nordeste, até 54,3 \%, na região Sul.
Há diferença estatisticamente significativa apenas entre as regiões Nordeste e Sul.
Considerando -se o desempenho dos alunos de Licenciatura, nota -se ainda que 43,9 \% consideraram a extensão da prova Adequada, no quarto de desempenho inferior, e 49,8 \%, no de melhor desempenho ( quarto superior ).
Nos quartos intermediários, essa proporção foi de 46,5 \%, no segundo quarto, e de 46,5 \%, no terceiro.
No Gráfico 4.10a, observa -se que as proporções de estudantes de Licenciatura que consideraram a prova Longa ou Muito longa apresentaram tendência decrescente em função dos quartos de desempenho, sendo a maior proporção dessa resposta no primeiro quarto de desempenho ( 46,2 \% ), e a menor, no quarto superior ( 41,7 \% ).
Pode -se constatar que não há diferença estatisticamente significativa entre os quartos de desempenho.
Considerando -se o desempenho dos alunos de Bacharelado, nota -se ainda que 46,0 \% consideraram a extensão da prova Adequada, no quarto de desempenho inferior, e 48,3 \%, no de melhor desempenho ( quarto superior ).
Nos quartos intermediários, essa proporção foi de 42,8 \%, no segundo quarto, e de 39,5 \%, no terceiro quarto.
No Gráfico 4.10b, pode -se constatar que não há diferença estatisticamente significativa entre as proporções de estudantes de Bacharelado que consideraram a prova Longa ou Muito longa em função dos quartos de desempenho, sendo a menor proporção dessa resposta no último quarto de desempenho ( 42,2 \% ), e a maior, no terceiro quarto ( 51,9 \% ).
Considerando -se a extensão da prova, o percentual de estudantes de Licenciatura que avaliaram a prova como Longa ou Muito longa, em relação ao tempo total destinado à sua resolução, foi maior entre os estudantes de Universidades, alcançando 44,5 \% ; entre os estudantes de Faculdades, a Organização Acadêmica com a menor incidência, esse índice atingiu 36,2 \% ; e entre os de Centros Universitários, 40,8 \%.
No Gráfico 4.11a, é possível observar que não há diferença estatisticamente significativa entre os tipos de Organização Acadêmica.
A proporção de presentes de Licenciatura à prova que consideraram a sua extensão Adequada, em relação ao tempo total destinado à sua resolução, esteve entre 44,7 \%, nas Universidades, e 53,1 \%, nas Faculdades ( ver também Tabela II.6a, no Anexo IIA ).
Considerando -se a extensão da prova, o percentual de estudantes de Bacharelado que avaliaram a prova como Longa ou Muito longa, em relação ao tempo total destinado à sua resolução, foi maior entre os estudantes de Centros Universitários, a proporção foi de 49,2 \% ; entre os de Faculdades, a Organização Acadêmica com a menor incidência, foi de 36,8 \%.
No Gráfico 4.11b, é possível observar que somente a diferença entre Universidades ( 47,6 \% ) e Faculdades é estatisticamente significativa.
Nas Organizações Acadêmicas, a proporção de presentes de Bacharelado à prova que consideraram a sua extensão como Adequada, em relação ao tempo total destinado à sua resolução esteve entre 39,9 \%, nos Centros Universitários, e 52,2 \%, nas Faculdades ( ver também Tabela II.6b, no Anexo IIB ).
O percentual de estudantes de Licenciatura que consideraram a prova Longa ou Muito longa, em relação ao tempo total destinado à sua resolução, foi menor entre os estudantes de cursos de IES Privadas ( 41,3 \% ), não havendo diferença estatisticamente significativa da proporção registrada em relação aos de IES Públicas ( 44,9 \% ).
Já a proporção de presentes de Licenciatura à prova que consideraram a extensão Adequada foi de 43,9 \%, entre os de IES Públicas, e de 49,9 \%, entre os de IES Privadas ( ver Gráfico 4.12a e Tabela II.6a, no Anexo IIA ).
O percentual de estudantes de Bacharelado que consideraram a prova como Longa ou Muito longa, em relação ao tempo total destinado à sua resolução, foi menor entre os estudantes de cursos de IES Privadas ( 43,4 \% ), sem diferença estatisticamente significativa da proporção registrada em IES Públicas ( 50,0 \% ).
Já a proporção de presentes de Bacharelado à prova que consideraram a extensão como Adequadas foi de 39,9 \%, entre os estudantes de IES Públicas, e de 47,0 \%, entre os de IES Privadas ( ver Gráfico 4.12b e Tabela II.6b, no Anexo IIB ). 4.3 COMPREENSÃO DOS ENUNCIADOS DAS QUESTÕES 4.3.1 COMPONENTE DE FORMAÇÃO GERAL Em relação aos enunciados das questões do Componente de Formação Geral ( Questão 4 ), as opiniões foram positivas, já que 75,3 \% dos alunos de Licenciatura avaliados consideraram todos ou a maioria dos enunciados das questões “ claros e objetivos ” ( Gráfico 4.13a, Gráfico 4.14a e, no Anexo IIA, a Tabela II.7a ).
Na análise regional, a percentagem de estudantes de Licenciatura que avaliaram que todos ou a maioria dos enunciados das questões do Componente de Formação Geral estavam “ claros e objetivos ” variou de 72,8 \%, na região Nordeste, a 77,8 \%, na região Norte.
A análise das percepções dos estudantes de Licenciatura sobre a clareza e a objetividade dos enunciados permite afirmar que todos, ou a maioria dos enunciados de questões relativas ao Componente de Formação Geral, foram considerados “ claros e objetivos ” para a maior parte dos respondentes ( maior ou igual a 72,8 \%, em todas as regiões, e maior ou igual a 62,1 \%, em todos os quartos de desempenho ).
Em relação aos enunciados das questões do Componente de Formação Geral ( Questão 4 ), as opiniões foram positivas, já que 76,1 \% dos alunos de Bacharelado avaliados consideraram todos ou a maioria dos enunciados das questões “ claros e objetivos ” ( Gráfico 4.13b, Gráfico 4.14b, e, no Anexo IIB, a Tabela II.7b ).
Na análise regional, a percentagem de estudantes de Bacharelado que avaliaram que todos ou a maioria dos enunciados das questões do Componente de Formação Geral estavam “ claros e objetivos ” variou de 66,7 \%, na região Norte, a 84,6 \%, na região Centro-Oeste.
A análise das percepções dos estudantes de Bacharelado sobre a clareza e objetividade dos enunciados permite afirmar que todos, ou a maioria dos enunciados de questões relativas ao Componente de Formação Geral, foram considerados “ claros e objetivos ” pela maior parte dos respondentes ( maior ou igual a 66,7 \%, em todas as regiões, e maior ou igual a 69,0 \%, em todos os quartos de desempenho ).
Segundo o desempenho, observa -se que a proporção dos que emitiram essa opinião cresce, conforme o desempenho aumenta, com diferenças estatisticamente significativas entre o quarto inferior de desempenho e os demais, assim como entre o segundo quarto e o quarto superior.
No quarto superior, a clareza e a objetividade de todos ou da maioria dos enunciados das questões foram percebidas por 82,7 \%.
Segundo o desempenho, observa -se que a proporção dos que emitiram essa opinião cresce conforme o desempenho aumenta até o terceiro quarto, com diferenças estatisticamente significativas entre o quarto inferior e os dois quartos superiores de desempenho.
No quarto superior, a clareza e a objetividade de todos ou da maioria dos enunciados das questões foram percebidas por 79,0 \%.
A proporção dos presentes de Licenciatura que consideraram que todos ou a maioria dos enunciados das questões do Componente de Formação Geral estavam “ claros e objetivos ” foi maior entre os estudantes de Universidades, chegando a 76,0 \% ; entre os de Faculdades, a Organização Acadêmica com a menor incidência, esse índice atingiu 71,3 \% ; e entre os de Centros Universitários, 73,9 \%.
No Gráfico 4.15a, é possível observar que as diferenças entre os tipos de Organização Acadêmica não são estatisticamente significativas ( ver também Tabela II.8a, no Anexo IIA ).
A proporção dos presentes de Bacharelado que consideraram que todos ou a maioria dos enunciados das questões do Componente de Formação Geral estavam “ claros e objetivos ” foi maior entre os estudantes de Centros Universitários, chegando a 78,6 \% ; entre os de Faculdades, a Organização Acadêmica com a menor incidência, esse índice atingiu 73,7 \% ; e entre os de Universidades, 76,0 \%.
No Gráfico 4.15b, é possível observar que não há diferença estatisticamente significativa entre as Organizações Acadêmicas ( ver também Tabela II.8b, no Anexo IIB ).
O percentual de estudantes de Licenciatura que consideraram que todos ou a maioria dos enunciados das questões do Componente de Formação Geral estavam “ claros e objetivos ” foi maior entre os estudantes de cursos de IES Públicas ( 75,4 \% ), mas sem diferença estatisticamente significativa da proporção registrada em IES Privadas ( 75,2 \% ).
Um maior detalhamento pode ser encontrado no Gráfico 4.16a e na Tabela II.8a, no Anexo IIA.
O percentual de estudantes de Bacharelado que consideraram que todos ou a maioria dos enunciados das questões do Componente de Formação Geral estavam “ claros e objetivos ” foi maior entre os estudantes de cursos de IES Privadas ( 76,6 \% ), sem diferença estatisticamente significativa da proporção registrada em IES Públicas ( 75,3 \% ).
Um maior detalhamento pode ser encontrado no Gráfico 4.16b e na Tabela II.8b, no Anexo IIB. 4.3.2 COMPONENTE DE CONHECIMENTO ESPECÍFICO Em relação aos enunciados das questões do Componente de Conhecimento Específico da prova, para 74,7 \% dos estudantes de Licenciatura avaliados da Área de Filosofia, a clareza e a objetividade ( Questão 5 ) estavam presentes em todas ou na maioria das questões ( Gráfico 4.17a a Gráfico 4.20a, e no Anexo IIA, as Tabelas II.9a e II.10a ).
A maioria dos estudantes de Licenciatura de todas as Grandes Regiões brasileiras considerou “ claros e objetivos ” todos ou a maioria dos enunciados das questões do Componente de Conhecimento Específico da prova, percentual sempre maior ou igual a 73,4 \%.
As diferenças entre as regiões não são estatisticamente significativas.
Em relação aos enunciados das questões do Componente de Conhecimento Específico da prova, para 71,3 \% dos estudantes de Bacharelado avaliados da Área de Filosofia, a clareza e a objetividade ( Questão 5 ) estavam presentes em todas ou na maioria das questões ( Gráfico 4.17b a Gráfico 4.20b, e no Anexo IIB, as Tabelas II.9b e II.10b ).
A maioria dos estudantes de Bacharelado de todas as Grandes Regiões brasileiras considerou “ claros e objetivos ” todos ou a maioria dos enunciados das questões do Componente de Conhecimento Específico da prova, percentual sempre maior ou igual a 59,5 \%.
As diferenças entre as regiões não são estatisticamente significativas.
A proporção de estudantes de Licenciatura que consideraram os enunciados das questões “ claros e objetivos ” apresenta uma tendência crescente em relação ao aumento de desempenho : mais elevada no quarto superior ( 82,0 \% ), se comparada ao quarto inferior de desempenho ( 62,3 \% ).
As diferenças entre o quarto inferior e os demais quartos de desempenho são estatisticamente significativas, assim como entre o segundo quarto e o último quarto.
A proporção de estudantes de Bacharelado que consideraram os enunciados das questões “ claros e objetivos ” apresenta uma tendência crescente em relação ao aumento de desempenho até o terceiro quarto : mais elevada no terceiro quarto ( 76,3 \% ), se comparada ao quarto inferior de desempenho ( 64,4 \% ).
Apenas a diferença entre o primeiro e o terceiro quartos de desempenho é estatisticamente significativa.
A proporção dos presentes de Licenciatura que consideraram que todos ou a maioria dos enunciados das questões do Componente de Conhecimento Específico estavam “ claros e objetivos ” foi maior entre os estudantes de Faculdades, atingindo 75,4 \%.
Entre os estudantes de Centro Universitários, a Organização Acadêmica com a menor incidência, essa proporção atingiu 73,7 \%, sendo intermediária a proporção entre os de Universidades ( 75,0 \% ).
No Gráfico 4.19a, é possível observar que as diferenças entre os tipos de Organização Acadêmica não são estatisticamente significativas ( ver também Tabela II.10a, no Anexo IIA ).
A proporção dos presentes de Bacharelado que consideraram que todos ou a maioria dos enunciados das questões do Componente de Conhecimento Específico estavam “ claros e objetivos ” foi maior entre os estudantes de Centro Universitários, atingindo 73,8 \%.
Entre os estudantes de Universidades, a Organização Acadêmica com a menor incidência, essa proporção atingiu 70,7 \%, sendo intermediária a proporção entre os de Faculdades ( 71,1 \% ).
No Gráfico 4.19b, é possível observar que não há diferença estatisticamente significativa entre as Organizações Acadêmicas ( ver também Tabela II.10b, no Anexo IIB ).
O percentual de estudantes de Licenciatura que consideraram que todos ou a maioria dos enunciados das questões do Componente de Conhecimento Específico estavam “ claros e objetivos ” foi maior entre os estudantes de cursos de IES Públicas ( 74,9 \% ), sem diferença estatisticamente significativa da proporção registrada em IES Privadas ( 74,5 \% ).
Ver Gráfico 4.20a e Tabela II.10a, no Anexo IIA, para maior detalhamento.
O percentual de estudantes de Bacharelado que consideraram que todos ou a maioria dos enunciados das questões do Componente de Conhecimento Específico estavam “ claros e objetivos ” foi maior entre os estudantes de cursos de IES Privadas ( 71,3 \% ), sem uma diferença estatisticamente significativa da proporção registrada em IES Públicas ( 71,2 \% ).
Ver Gráfico 4.20b e Tabela II.10b, no Anexo IIB, para um maior detalhamento. 4.4 SUFICIÊNCIA DAS INFORMAÇÕES/INSTRUÇÕES FORNECIDAS Ao se avaliarem as informações/instruções fornecidas para a resolução das questões ( Questão 6 ), constatou -se que 82,3 \% dos respondentes de Licenciatura da Área de Filosofia de todo o Brasil as consideraram até excessivas ou suficientes em todas ou na maioria das questões ( Gráfico 4.21a, Gráfico 4.22a e, no Anexo IIA, a Tabela II.11a ).
Quanto à distribuição de respondentes pelas Grandes Regiões, observa -se que a proporção de estudantes que consideraram as informações/instruções fornecidas até excessivas ou suficientes em todas ou na maioria das questões foi sempre superior ou igual a 78,0 \%, chegando a 84,0 \%, na região Sul.
Não existem diferenças estatisticamente significativas entre as regiões.
Ao avaliarem as informações/instruções fornecidas para a resolução das questões ( Questão 6 ), 85,1 \% dos respondentes de Bacharelado da Área de Filosofia de todo o Brasil as consideraram até excessivas ou suficientes em todas ou na maioria das questões ( Gráfico 4.21b, Gráfico 4.22b, e, no Anexo IIB, a Tabela II.11b ).
Quanto à distribuição de respondentes pelas Grandes Regiões, observa -se que a proporção de estudantes que consideraram as informações/instruções fornecidas até excessivas ou suficientes em todas ou na maioria das questões foi sempre superior ou igual a 82,0 \%, chegando a 89,0 \%, na região Centro-Oeste.
Não existem diferenças estatisticamente significativas entre as regiões.
Levando -se em conta o desempenho dos participantes de Licenciatura, observa -se uma tendência crescente das proporções de participantes que avaliaram as informações/instruções como até excessivas ou suficientes em todas ou na maioria das questões, havendo diferenças estatisticamente significativas entre o primeiro quarto e os demais quartos de desempenho, conforme mostrado no Gráfico 4.22a.
O percentual foi mais elevado no quarto superior ( 87,8 \% ), percentual superior à média nacional ( 82,3 \% ).
Já no quarto inferior, a suficiência das informações/instruções declarada como até excessiva, em todas ou na maioria das questões foi percebida por 71,2 \% dos respondentes.
Levando -se em conta o desempenho dos participantes de Bacharelado, observa -se uma tendência crescente das proporções de participantes que avaliaram as informações/instruções como até excessivas ou suficientes em todas ou na maioria das questões, havendo diferença estatisticamente significativa apenas entre o quarto inferior e o quarto superior de desempenho, como mostrado no Gráfico 4.22b.
O percentual foi mais elevado no quarto superior ( 88,1 \% ), percentual superior à média nacional ( 85,1 \% ).
Já no quarto inferior, as informações/instruções consideradas como até excessivas, em todas ou na maioria das questões foram percebidas por 80,2 \% dos respondentes.
Considerando -se as informações/instruções fornecidas para a resolução das questões ( Questão 6 ), a proporção que as considerou até excessivas ou suficientes em todas ou na maioria das questões foi maior entre os estudantes de Licenciatura de Centros Universitários, totalizando 84,6 \%.
As Universidades foram a Organização Acadêmica em que houve a menor proporção : 81,4 \%.
No Gráfico 4.23a, é possível observar que as diferenças entre as Organizações Acadêmicas não são estatisticamente significativas ( ver também Tabela II.12a, no Anexo IIA ).
Considerando -se as informações/instruções fornecidas para a resolução das questões ( Questão 6 ), a proporção que as consideraram até excessivas ou suficientes em todas ou na maioria das questões foi maior entre os estudantes de Bacharelado de Centros Universitários, uma proporção de 86,9 \%, enquanto entre os de Universidades, a Organização Acadêmica com a menor incidência, a proporção atingiu 84,2 \%.
No Gráfico 4.23b, é possível observar que não existe diferença estatisticamente significativa entre as Organizações Acadêmicas ( ver também Tabela II.12b, no Anexo IIB ).
O percentual de estudantes de Licenciatura que consideraram que as informações/instruções fornecidas para a resolução das questões eram até excessivas ou suficientes em todas ou na maioria das questões foi menor entre os estudantes de cursos de IES Públicas ( 80,8 \% ), não havendo diferença estatisticamente significativa da proporção registrada em IES Privadas ( 84,0 \% ) ( ver Gráfico 4.24a e Tabela II.12a, no Anexo IIA ).
O percentual de estudantes de Bacharelado que consideraram que as informações/instruções fornecidas para a resolução das questões eram até excessivas ou suficientes em todas ou na maioria das questões foi menor entre os estudantes de cursos de IES Públicas ( 84,5 \% ), sem uma diferença estatisticamente significativa da proporção registrada em IES Privadas ( 85,6 \% ) ( ver Gráfico 4.24b e Tabela II.12b, no Anexo IIB ). 4.5 DIFICULDADE ENCONTRADA AO RESPONDER À PROVA Indagados sobre as dificuldades com as quais se depararam ao responder à prova ( Questão 7 ), 15,7 \% dos estudantes de Licenciatura apontaram o Desconhecimento do conteúdo.
Para 43,2 \%, a Forma diferente de abordagem do conteúdo foi indicada como dificuldade.
Já a Falta de motivação para fazer a prova foi a dificuldade apontada por 17,3 \% dos respondentes.
Considerando -se todo o Brasil, 16,5 \% dos respondentes afirmaram que não tiveram qualquer tipo de dificuldade para responder à prova ( Tabelas II.13a e II.14a, no Anexo IIA ).
Nos Gráficos 4.25a a 4.28a, são apresentados os percentuais de estudantes de Licenciatura que apontaram o Desconhecimento do conteúdo como “ dificuldade ao responder à prova ”.
Na análise por Grandes Regiões, o percentual de inscritos e presentes de Licenciatura que apontaram o Desconhecimento do conteúdo como “ dificuldade ao responder à prova ” não superou 18,3 \%.
Os percentuais variaram de 12,6 \%, na região Centro-Oeste, a 18,3 \%, na região Nordeste.
A Forma diferente de abordagem do conteúdo foi a escolha modal dos estudantes de Licenciatura, com percentuais que variaram de 42,2 \% ( região Nordeste ) a 49,2 \% ( região Centro-Oeste ).
O percentual de alunos que citaram a Falta de motivação como dificuldade variou de 13,8 \% ( região Norte ) a 17,9 \% ( região Sul ).
Os que declararam não ter qualquer tipo de dificuldade para responder à prova variaram de 12,9 \%, na região Norte, a 17,8 \%, na região Centro-Oeste.
Indagados sobre as dificuldades com as quais se depararam ao responder à prova ( Questão 7 ), 17,9 \% dos estudantes de Bacharelado apontaram o Desconhecimento do conteúdo.
Para 44,1 \%, a Forma diferente de abordagem do conteúdo foi indicada como dificuldade.
Já a Falta de motivação para fazer a prova foi a dificuldade apontada por 16,2 \% dos respondentes.
Considerando -se todo o Brasil, 15,2 \% dos respondentes afirmaram que não tiveram qualquer tipo de dificuldade para responder à prova ( Tabelas II.13b e II.14b, no Anexo IIB ).
Nos Gráficos 4.25b a 4.28b, são apresentados os percentuais de estudantes de Bacharelado que apontaram o Desconhecimento do conteúdo como “ dificuldade ao responder à prova ”.
Na análise por Grandes Regiões, o percentual de inscritos e presentes de Bacharelado que apontaram o Desconhecimento do conteúdo como “ dificuldade ao responder à prova ” não superou 40,0 \%.
Os percentuais variaram de 15,4 \%, na região Centro-Oeste a 40,0 \%, na região Norte.
A Forma diferente de abordagem do conteúdo foi a escolha modal dos estudantes de Bacharelado, com percentuais que variaram de 37,1 \% ( região Norte ) a 47,1 \% ( região Nordeste ).
O percentual de alunos que citaram a Falta de motivação como dificuldade variou de 5,7 \% ( região Norte ) a 16,7 \% ( região Sul ).
Os que declararam não ter qualquer tipo de dificuldade para responder à prova variaram de 8,6 \%, na região Norte, a 17,6 \%, na região Centro-Oeste.
Em relação aos quartos de desempenho, o Desconhecimento do conteúdo foi a opção escolhida por 15,2 \% dos estudantes de Licenciatura do quarto inferior e por 15,0 \% do quarto superior.
Não são observadas diferenças estatisticamente significativas entre os quartos de desempenho.
A alternativa modal entre os alunos de Licenciatura, quando agregados pelos quartos de desempenho, para a dificuldade encontrada, foi causada pela Forma diferente de abordagem do conteúdo : 46,3 \%, no primeiro quarto, e 39,0 \%, no último quarto.
Em relação aos quartos de desempenho, o Desconhecimento do conteúdo foi a opção escolhida por 20,1 \% dos estudantes de Bacharelado do quarto inferior e por 16,0 \% do quarto superior.
Não são observadas diferenças estatisticamente significativas entre os quartos de desempenho.
A alternativa modal entre os alunos de Bacharelado, quando agregados pelos quartos de desempenho, para a dificuldade encontrada, foi causada pela Forma diferente de abordagem do conteúdo : variando de 40,1 \%, no quarto superior, a 49,6 \%, no segundo quarto.
Na análise por Organização Acadêmica, o percentual de inscritos e presentes de Licenciatura que apontaram o Desconhecimento do conteúdo como “ dificuldade ao responder à prova ” não superou 19,7 \%.
Os percentuais variaram de 14,0 \%, nos Centros Universitários, a 19,7 \%, nas Faculdades.
A Forma diferente de abordagem do conteúdo foi a escolha modal dos estudantes de Licenciatura, com percentuais que variaram de 41,8 \% ( Universidades ) a 52,0 \% ( Faculdades ).
O percentual de alunos de Licenciatura que citaram a Falta de motivação como dificuldade variou de 7,1 \% ( Faculdades ) a 18,2 \% ( Universidades ).
Os que declararam não ter qualquer tipo de dificuldade para responder à prova variaram de 15,7 \%, nas Faculdades, a 17,7 \%, nos Centros Universitários.
No Gráfico 4.27a, é possível observar que não há diferenças estatisticamente significativas entre Faculdades, Centros Universitários e Universidades ( ver também Tabela II.14a, no Anexo IIA ).
Na análise por Organização Acadêmica, o percentual de inscritos e presentes de Bacharelado que apontaram o Desconhecimento do conteúdo como “ dificuldade ao responder à prova ” não superou 22,6 \%.
Os percentuais variaram de 15,2 \%, entre os de Centros Universitários a 22,6 \%, entre os de Faculdades.
A Forma diferente de abordagem do conteúdo foi a escolha modal dos estudantes de Bacharelado, com percentuais que variaram de 41,7 \% ( Universidades ) a 48,7 \% ( Faculdades ).
O percentual de alunos de Bacharelado que citaram a Falta de motivação como dificuldade variou de 13,5 \% ( Centros Universitários ) a 17,4 \% ( Universidades ).
Os que declararam não ter qualquer tipo de dificuldade para responder à prova variaram de 10,6 \%, nas Faculdades, a 17,7 \%, nos Centros Universitários.
No Gráfico 4.27b, é possível observar não há diferenças estatisticamente significativas entre Faculdades, Centros Universitários e Universidades ( ver também Tabela II.14b, no Anexo IIB ).
Em relação às Categorias Administrativas, o Desconhecimento do conteúdo foi a opção escolhida por 17,3 \% dos estudantes de Licenciatura, nas IES Públicas, e por 13,9 \%, nas IES Privadas.
Não são observadas diferenças estatisticamente significativas entre essas duas categorias.
A alternativa modal entre os alunos de Licenciatura, quando agregados pelas Categorias Administrativas, para a dificuldade encontrada, foi causada pela Forma diferente de abordagem do conteúdo : 41,6 \%, entre os estudantes de IES Públicas, e 45,1 \%, entre os de IES Privadas ( ver Gráfico 4.28a e Tabela II.14a, no Anexo IIA ).
Em relação às Categorias Administrativas, o Desconhecimento do conteúdo foi a opção escolhida por 18,7 \% dos estudantes de Bacharelado, nas IES Públicas, e por 17,3 \%, nas IES Privadas.
Não são observadas diferenças estatisticamente significativas entre essas duas categorias.
A alternativa modal entre os alunos de Bacharelado, quando agregados pelas Categorias Administrativas, para a dificuldade encontrada, foi causada pela Forma diferente de abordagem do conteúdo : 39,8 \%, entre os estudantes de IES Públicas, e 46,9 \%, entre os de IES Privadas ( ver Gráfico 4.28b e Tabela II.14b, no Anexo IIB ). 4.6 CONTEÚDOS DAS QUESTÕES OBJETIVAS DA PROVA Ao se analisarem os conteúdos das questões objetivas da prova ( Questão 8 ), verifica -se que um percentual pequeno dos estudantes de Licenciatura avaliados, apenas 6,6 \%, afirmou que não estudou ainda a maioria desses conteúdos ( Gráficos 4.29a a Gráfico 4.32a, e as Tabelas II.15a e II.16a, no Anexo IIA ).
A maioria ( 67,2 \% ) afirmou ter estudado e aprendido muitos ou todos os conteúdos avaliados.
Na análise por Grande Região, a proporção de respondentes de Licenciatura que escolheram a opção não estudou ainda a maioria desses conteúdos foi pequena.
Observa -se que, nas regiões Norte ( 7,6 \% ) e Sul ( 7,2 \% ), as proporções foram maiores do que a média nacional ( 6,6 \% ).
Nas regiões Nordeste, Sudeste e Sul, a situação foi inversa, com proporções mais baixas, correspondendo a, respectivamente, 6,3 \%, 6,2 \% e 6,3 \%.
Não são observadas diferenças estatisticamente significativas entre as regiões.
Em todas as regiões, a maioria dos presentes afirmou ter estudado e aprendido muitos ou todos os conteúdos, com proporções variando entre 59,3 \%, na região Norte, a 70,4 \%, na região Centro-Oeste.
Ao analisar os conteúdos das questões objetivas da prova ( Questão 8 ), um percentual pequeno dos estudantes de Bacharelado avaliados, apenas 7,2 \%, afirmaram que não estudou ainda a maioria desses conteúdos ( Gráficos 4.29b a Gráfico 4.32b, e as Tabelas II.15b e II.16b, no Anexo IIB ).
A maioria ( 65,2 \% ) afirmou ter estudado e aprendido muitos ou todos os conteúdos avaliados.
Na análise por Grande Região, a proporção de respondentes de Bacharelado que escolheram a opção não estudou ainda a maioria desses conteúdos foi pequena.
Observa -se que, nas regiões Norte ( 20,0 \% ), Nordeste ( 7,8 \% ) e Sudeste ( 8,4 \% ), as proporções foram maiores do que a média nacional ( 7,2 \% ).
Nas regiões Sul ( 4,1 \% ) e Centro-Oeste ( 4,4 \% ), a situação foi inversa : uma proporção mais baixa.
Não são observadas diferenças estatisticamente significativas entre as regiões.
Em todas as regiões, a maioria dos presentes afirmou ter estudado e aprendido muitos ou todos os conteúdos, com proporções variando entre 31,4 \%, na região Norte, a 70,7 \%, na região Sul.
Considerando -se separadamente as opiniões de estudantes de Licenciatura dos quatro quartos de desempenho, observa -se que, no quarto inferior, 12,6 \% ofereceram como resposta que não estudou ainda a maioria desses conteúdos, sendo 1,8 \% os do quarto superior com a mesma resposta.
A diferença entre os alunos que optaram por esse motivo de dificuldade no quarto inferior em relação aos demais quartos de desempenho é estatisticamente significativa, bem como a diferença entre o quarto superior em relação aos demais quartos de desempenho.
Tendo -se em conta o quarto superior, 86,1 \% dos alunos de Licenciatura afirmaram ter estudado e aprendido muitos ou todos os conteúdos.
No outro extremo, no primeiro quarto, 45,1 \% optaram pelas mesmas categorias.
As proporções são crescentes com o desempenho.
Considerando -se, separadamente, as opiniões de estudantes de Bacharelado dos quatro quartos de desempenho, observa -se que, no quarto inferior, 15,9 \% ofereceram como resposta que não estudou ainda a maioria desses conteúdos, sendo 3,5 \% os dos dois quartos superiores com a mesma resposta.
A diferença entre os alunos que optaram por esse motivo de dificuldade no quarto inferior, em relação aos demais quartos de desempenho, é estatisticamente significativa.
Tendo -se em conta o quarto superior, 80,4 \% dos alunos de Bacharelado afirmaram ter estudado e aprendido muitos ou todos os conteúdos.
No outro extremo, no primeiro quarto, 44,3 \% optaram pelas mesmas categorias.
As proporções são crescentes com o desempenho.
Na análise por Organização Acadêmica, a proporção de respondentes de Licenciatura que escolheram a opção não estudou ainda a maioria desses conteúdos foi pequena.
Observa -se que as proporções nas Universidades ( 5,9 \% ) e Faculdades ( 3,8 \% ) foram menores do que a média nacional ( 6,6 \% ).
Nos Centros Universitários, a situação foi inversa : a proporção foi mais alta ( 8,7 \% ).
No Gráfico 4.31a, é possível observar que as diferenças entre os tipos de Organização Acadêmica não são estatisticamente significativas.
Em todos os tipos de Organização Acadêmica, a maioria dos presentes afirmou ter estudado e aprendido muitos ou todos os conteúdos, com proporções variando entre 62,3 \%, nas Faculdades, a 68,2 \%, nas Universidades ( ver também Tabela II.16a, no Anexo IIA ).
Na análise por Organização Acadêmica, a proporção de respondentes de Bacharelado que escolheram a opção não estudou ainda a maioria desses conteúdos foi pequena.
Observa -se que, nas Universidades ( 7,4 \% ) e Centros Universitários ( 7,6 \% ) as proporções foram maiores do que a média nacional ( 7,2 \% ).
Nas Faculdades ( 6,2 \% ), a situação foi inversa : a proporção foi mais baixa.
No Gráfico 4.31b, é possível observar que a diferença entre as Organizações Acadêmicas não é estatisticamente significativa ( ver também Tabela II.16b, no Anexo IIB ).
Em todos os tipos de Organização Acadêmica, a maioria dos presentes afirmou ter estudado e aprendido muitos ou todos os conteúdos, com proporções variando entre 54,6 \%, nas Faculdades, a 68,6 \%, nas Universidades.
Em relação às Categorias Administrativas, observa -se que, nas IES Públicas, 5,9 \% ofereceram como resposta que não estudou ainda a maioria desses conteúdos, sendo 7,2 \% nas IES Privadas com a mesma resposta.
A diferença entre os alunos de Licenciatura que optaram por esse motivo de dificuldade nas duas categorias não é estatisticamente significativa.
Considerando -se as IES Públicas, 67,8 \% dos alunos de Licenciatura afirmaram ter estudado e aprendido muitos ou todos os conteúdos, sendo de 66,4 \% a proporção dos que escolheram essa resposta nas IES Privadas ( ver Gráfico 4.32a e Tabela II.16a, no Anexo IIA ).
Em relação às Categorias Administrativas, observa -se que, nas IES Públicas, 8,5 \% ofereceram como resposta que não estudou ainda a maioria desses conteúdos, sendo 6,4 \% nas IES Privadas com a mesma resposta.
A diferença entre os alunos de Bacharelado que optaram por esse motivo de dificuldade nas duas categorias não é estatisticamente significativa.
Considerando -se as IES Públicas, 71,8 \% dos alunos de Bacharelado afirmaram ter estudado e aprendido muitos ou todos os conteúdos.
Na outra categoria, 60,9 \% optaram pelas mesmas categorias ( ver Gráfico 4.32b e Tabela II.16b, no Anexo IIB ). 4.7 TEMPO GASTO PARA CONCLUIR A PROVA Ao responderem sobre o tempo de conclusão da prova ( Questão 9 ), mais de três quartos dos estudantes de Licenciatura ( 75,3 \% ) afirmaram ter gasto Entre duas e quatro horas ( Gráfico 4.33a a Gráfico 4.36a e, no Anexo IIA, as Tabelas II.17a e II.18a ).
Considerando -se as cinco Grandes Regiões brasileiras, os percentuais dos que utilizaram Entre duas e quatro horas para finalizar a prova nas regiões Norte ( 72,4 \% ), Sul ( 73,4 \% ) e Centro-Oeste ( 74,7 \% ) foram inferiores ao percentual nacional.
Nas demais Grandes Regiões, os percentuais de alunos de Licenciatura que dispensaram Entre duas e quatro horas para concluir a prova ficaram acima : 76,5 \% na região Nordeste e 76,5 \% na região Sudeste, como mostrado no Gráfico 4.33a.
Não há diferenças estatisticamente significativas entre as regiões.
Ao responderem sobre o tempo de conclusão da prova ( Questão 9 ), menos de três quartos dos estudantes de Bacharelado ( 73,9 \% ) afirmaram ter gasto Entre duas e quatro horas ( Gráfico 4.33b a Gráfico 4.36b e, no Anexo IIB, as Tabelas II.17b e II.18b ).
Considerando -se as cinco Grandes Regiões brasileiras, os percentuais dos que utilizaram Entre duas e quatro horas para finalizar a prova nas regiões Norte ( 62,9 \% ) e Nordeste ( 70,1 \% ) foram inferiores ao percentual nacional.
Nas demais Grandes Regiões, os percentuais de alunos de Bacharelado que dispensaram Entre duas e quatro horas para concluir a prova foram : 74,9 \%, na região Sudeste ; 76,0 \%, na região Sul ; e 74,3 \%, na região Centro-Oeste, como mostrado no Gráfico 4.33b.
A diferença entre as regiões não é estatisticamente significativa.
Uma vez analisadas as alternativas escolhidas pelos estudantes de Licenciatura que se situam nos diferentes quartos de desempenho, observa -se uma tendência crescente : uma proporção maior de participantes no quarto superior declarou ter gasto Entre duas e quatro horas para concluir a prova quando comparada com a proporção de participantes nos quartos inferiores.
Há diferenças estatisticamente significativas entre o quarto inferior e os dois quartos superiores de desempenho, bem como entre o segundo e o quarto superior de desempenho.
Uma vez analisadas as alternativas escolhidas pelos estudantes de Bacharelado que se situam nos diferentes quartos de desempenho, observa -se uma tendência crescente : uma proporção maior de participantes no quarto superior declarou ter gasto Entre duas e quatro horas para concluir a prova quando comparada com a proporção de participantes nos quartos inferiores.
Há diferenças estatisticamente significativas entre os dois quartos inferiores e os dois quartos superiores de desempenho, mas não entre si.
Na análise por Organização Acadêmica, os percentuais dos que utilizaram Entre duas e quatro horas para finalizar a prova nas Faculdades ( 83,5 \% ) e Centros Universitários ( 75,5 \% ) foram superiores ao percentual nacional ( 75,3 \% ).
Nas Universidades, tal percentual foi de 74,8 \%, e, como mostrado no Gráfico 4.35a, há diferença estatisticamente significativa apenas entre as proporções das Universidades e Faculdades.
Na análise por Organização Acadêmica, o percentual dos que utilizaram Entre duas e quatro horas para finalizar a prova nas Faculdades ( 77,5 \% ) foi superior ao percentual nacional ( 73,8 \% ), enquanto nas Universidades e nos Centros Universitários, tais percentuais foram de 73,0 \% e 73,8 \%, respectivamente, como mostrado no Gráfico 4.35b.
As diferenças entre as Organizações Acadêmicas não são estatisticamente significativas.
Em relação às Categorias Administrativas, observa -se que, nas IES Públicas, 75,2 \% declararam ter gasto Entre duas e quatro horas para concluir a prova, sendo 75,5 \% nas IES Privadas com a mesma resposta.
A diferença entre os alunos de Licenciatura que optaram por esse motivo de dificuldade nas duas categorias não é estatisticamente significativa ( ver Gráfico 4.36a e Tabela II.18a, no Anexo IIA ).
Em relação às Categorias Administrativas, observa -se que, nas IES Públicas, 71,9 \% declararam ter gasto Entre duas e quatro horas para concluir a prova, sendo 75,2 \% nas IES Privadas com a mesma resposta.
A diferença entre os alunos de Bacharelado que optaram por esse motivo de dificuldade nas duas categorias não é estatisticamente significativa ( ver Gráfico 4.36b e Tabela II.18b, no Anexo IIB ).

% matched lemmas: 
\end{textsample}
